\documentclass[11pt]{article}

\usepackage[a4paper,margin=1in]{geometry}
\usepackage{amsmath,amssymb,bm}
\usepackage[T1]{fontenc}

\title{Implemented Positional Encoding in Training\_new}
\author{}
\date{}

\begin{document}
\maketitle

\paragraph{Notation.}
Let $G=(\mathcal{V},\mathcal{E})$ be the node graph used for PE.
For node $i\in\mathcal{V}$:
\[
\begin{aligned}
&d_i &&\text{: node degree},\\
&s_i\in\{0,1\} &&\text{: slack indicator},\\
&g_i\in\{0,1\} &&\text{: ground indicator},\\
&\mathcal{S}=\{u\in\mathcal{V}: s_u=1\} &&\text{: set of slack nodes.}
\end{aligned}
\]

\paragraph{BFS distance.}
BFS (Breadth-First Search) gives shortest path length in \emph{number of edges}.
Thus
\[
h_i = \min_{u\in\mathcal{S}} \operatorname{dist}_{\mathrm{BFS}}(i,u),
\qquad
h_i \leftarrow \min(h_i, H),
\]
where $H=\texttt{DK\_PE\_HOPCAP}$.

The node positional encoding is
\[
\bm{p}_i =
\begin{bmatrix}
\log(1+d_i) \\
\dfrac{h_i}{\max(10, H/3)} \\
s_i \\
g_i \\
\dfrac{d_i}{d_{\max}} \\
\dfrac{\log(1+d_i)}{\log(1+d_{\max})+10^{-6}}
\end{bmatrix}
\in \mathbb{R}^{6},
\qquad
d_{\max}=\max_{j\in\mathcal{V}} d_j.
\]

\paragraph{Why log terms.}
The transforms $\log(1+d_i)$ and $\log(1+d_{\max})$ compress heavy-tailed degree scale,
so high-degree nodes do not dominate the feature magnitude. Using $\log(1+x)$ is also
well-defined at $x=0$ and keeps features numerically stable.

Hence,
\[
\texttt{PE\_DIM}=6.
\]

Per-node encoder input is
\[
\bm{x}^{\mathrm{node}}_i
=
\big[\bm{V}^{\mathrm{init}}_i,\; \Delta\bm{V}_i\odot m_i,\; m_i,\; \bm{p}_i\big],
\]
where
\[
\begin{aligned}
&\bm{V}^{\mathrm{init}}_i\in\mathbb{R}^2 &&\text{: initial voltage (real, imag)},\\
&\Delta\bm{V}_i\in\mathbb{R}^2 &&\text{: voltage residual target},\\
&m_i\in\{0,1\} &&\text{: voltage visibility mask},\\
&\odot &&\text{: elementwise product},\\
&[\cdot,\cdot,\cdot] &&\text{: feature concatenation.}
\end{aligned}
\]

No-PE ablation:
\[
\bm{p}_i \leftarrow \bm{0}.
\]

\end{document}

